\documentclass[11pt]{article}

\usepackage[margin=1in]{geometry}
\usepackage{amsmath,amssymb}
\usepackage{booktabs}
\usepackage{array}
\usepackage{enumitem}
\usepackage{siunitx}
\usepackage{xcolor}
\usepackage{hyperref}

\hypersetup{
    colorlinks=true,
    linkcolor=blue!55!black,
    urlcolor=blue!55!black,
    citecolor=blue!55!black
}

\sisetup{per-mode=symbol}
\DeclareSIUnit\gpm{gpm}

\title{OASIS: Mathematical Model for a\\Liquid-Cooled GPU Data Center}
\author{Aarav Shandilya and the OASIS Engineering Team}
\date{\today}

\begin{document}

\maketitle

\begin{abstract}
OASIS is a coordinated architectural and engineering model of a high-density GPU
data center. The Autodesk Revit model defines the physical arrangement of racks,
coolant distribution units, pumps, piping, electrical equipment, and heat-rejection
equipment. The mathematical model evaluates whether the proposed cooling system can
maintain safe GPU and coolant temperatures during changing computational loads,
high outdoor temperatures, and equipment failures. This report defines the governing
energy, heat-transfer, fluid-flow, and numerical equations used in the model.
\end{abstract}

\tableofcontents
\newpage

\section{Project Objective}

The primary engineering question is:

\begin{quote}
How can a closed-loop liquid-cooling system be scaled across multiple GPU racks
while maintaining safe equipment temperatures and minimizing cooling energy and
water consumption?
\end{quote}

The cooling path is modeled as

\begin{center}
GPU $\rightarrow$ cold plate $\rightarrow$ secondary coolant loop
$\rightarrow$ coolant distribution unit (CDU) $\rightarrow$ facility loop
$\rightarrow$ dry cooler or chiller $\rightarrow$ outdoor environment.
\end{center}

The initial model uses lumped thermal nodes. Each node is treated as having a
single representative temperature. This approximation makes the equations suitable
for an RK4 differential-equation solver while preserving the dominant transient
thermal behavior.

\section{Symbols and Units}

\begin{table}[h!]
\centering
\small
\begin{tabular}{>{\(}l<{\)} p{8.3cm} l}
\toprule
\text{Symbol} & \text{Definition} & \text{SI unit} \\
\midrule
T_g & Effective GPU and cold-plate temperature & \si{\degreeCelsius} \\
T_s & Secondary coolant-loop temperature & \si{\degreeCelsius} \\
T_p & Primary facility-loop temperature & \si{\degreeCelsius} \\
T_a & Ambient or outdoor heat-sink temperature & \si{\degreeCelsius} \\
P_{\mathrm{IT}} & Total electrical power converted to IT heat & \si{\watt} \\
\dot Q & Heat-transfer rate & \si{\watt} \\
C & Lumped thermal capacitance & \si{\joule\per\kelvin} \\
m & Mass & \si{\kilogram} \\
c_p & Specific heat capacity at constant pressure & \si{\joule\per\kilogram\per\kelvin} \\
h & Convective heat-transfer coefficient & \si{\watt\per\square\metre\per\kelvin} \\
U & Overall heat-transfer coefficient & \si{\watt\per\square\metre\per\kelvin} \\
A & Effective heat-transfer area & \si{\square\metre} \\
G & Effective thermal conductance, equal to \(UA\) or \(hA\) & \si{\watt\per\kelvin} \\
R_{\mathrm{th}} & Thermal resistance, equal to \(1/G\) & \si{\kelvin\per\watt} \\
\dot m & Coolant mass-flow rate & \si{\kilogram\per\second} \\
\dot V & Coolant volumetric-flow rate & \si{\cubic\metre\per\second} \\
\rho & Coolant density & \si{\kilogram\per\cubic\metre} \\
\Delta p & Pressure difference or pressure loss & \si{\pascal} \\
\eta & Pump or equipment efficiency & dimensionless \\
\bottomrule
\end{tabular}
\end{table}

Temperature differences in degrees Celsius and kelvin have the same numerical value.
Absolute temperatures must be converted to kelvin when used in radiation equations.

\section{Electrical Load Converted to Heat}

Almost all electrical energy consumed by IT equipment eventually becomes heat inside
or near the data center. The rack heat load is approximated by

\begin{equation}
P_{\mathrm{rack}}(t)
= N_g P_{g,\mathrm{rated}}u_g(t)
+ P_{\mathrm{CPU}}(t)
+ P_{\mathrm{memory}}(t)
+ P_{\mathrm{network}}(t)
+ P_{\mathrm{other}}(t),
\end{equation}

where \(u_g(t)\) is the GPU utilization fraction. The full data-hall load is

\begin{equation}
P_{\mathrm{IT}}(t)=\sum_{r=1}^{N_r}P_{\mathrm{rack},r}(t).
\end{equation}

If identical racks are used with a load-diversity factor \(f_d\),

\begin{equation}
P_{\mathrm{IT,design}}=N_r f_d P_{\mathrm{rack,rated}}.
\end{equation}

\section{First Law of Thermodynamics}

For a control volume with no mechanical work other than separately modeled pump work,
the energy balance is

\begin{equation}
\frac{dE}{dt}=\dot Q_{\mathrm{in}}-\dot Q_{\mathrm{out}}+P_{\mathrm{generated}}.
\end{equation}

For a lumped node with approximately constant thermal capacitance,

\begin{equation}
E=CT,
\qquad
C=mc_p,
\end{equation}

so that

\begin{equation}
C\frac{dT}{dt}=\dot Q_{\mathrm{in}}-\dot Q_{\mathrm{out}}+P_{\mathrm{generated}}.
\end{equation}

This equation is the foundation of every thermal node in the OASIS model.

\section{Newton's Law of Cooling}

Convective heat transfer between a surface and a fluid is modeled using

\begin{equation}
\dot Q=hA(T_{\mathrm{hot}}-T_{\mathrm{cold}}).
\end{equation}

Because \(h\) and \(A\) appear together, the model defines an effective thermal
conductance

\begin{equation}
G=hA=\frac{1}{R_{\mathrm{th}}}.
\end{equation}

The heat transferred from the GPU assembly to the secondary coolant is therefore

\begin{equation}
\dot Q_{g\rightarrow s}=G_{gs}(T_g-T_s).
\end{equation}

\subsection{Single-Node Baseline Model}

Before solving the complete coupled model, a single-node model will be used for
verification:

\begin{equation}
C\frac{dT}{dt}=P-G(T-T_a).
\label{eq:one-node}
\end{equation}

For constant \(P\), \(G\), and \(T_a\), the analytical solution is

\begin{equation}
T(t)=T_a+\frac{P}{G}
+\left[T(0)-T_a-\frac{P}{G}\right]
e^{-Gt/C}.
\label{eq:analytical}
\end{equation}

The steady-state temperature is

\begin{equation}
T_{\infty}=T_a+\frac{P}{G},
\end{equation}

and the thermal time constant is

\begin{equation}
\tau=\frac{C}{G}.
\end{equation}

Equation~\eqref{eq:analytical} provides a direct test of the RK4 implementation.

\section{Coupled GPU and Coolant Model}

\subsection{GPU Node}

The GPU and cold-plate energy balance is

\begin{equation}
C_g\frac{dT_g}{dt}
=P_{\mathrm{IT}}(t)-G_{gs}(T_g-T_s).
\label{eq:gpu}
\end{equation}

Thus,

\begin{equation}
\frac{dT_g}{dt}
=\frac{P_{\mathrm{IT}}(t)-G_{gs}(T_g-T_s)}{C_g}.
\end{equation}

\subsection{Secondary Coolant Node}

The secondary coolant receives heat from the GPUs and rejects it through the CDU
heat exchanger:

\begin{equation}
C_s\frac{dT_s}{dt}
=G_{gs}(T_g-T_s)-G_{sp}(T_s-T_p).
\label{eq:secondary}
\end{equation}

\subsection{Primary Facility-Loop Node}

The primary loop receives heat through the CDU and rejects it to the outdoor system:

\begin{equation}
C_p\frac{dT_p}{dt}
=G_{sp}(T_s-T_p)-\dot Q_{\mathrm{reject}}.
\label{eq:primary}
\end{equation}

The unconstrained outdoor heat rejection is

\begin{equation}
\dot Q_{\mathrm{UA}}
=G_{pa}(T_p-T_a).
\end{equation}

To represent a finite equipment rating and prevent nonphysical reverse heat rejection,

\begin{equation}
\dot Q_{\mathrm{reject}}
=\min\left\{
\max\left[G_{pa}(T_p-T_a),0\right],
\dot Q_{\max}
\right\}.
\label{eq:qreject}
\end{equation}

Equations~\eqref{eq:gpu}--\eqref{eq:qreject} form the principal transient model.

\section{Coolant Heat Transport and Flow Sizing}

The sensible heat carried by a flowing coolant is

\begin{equation}
\dot Q=\dot m c_p(T_{\mathrm{return}}-T_{\mathrm{supply}}).
\label{eq:flowheat}
\end{equation}

The required mass-flow rate is

\begin{equation}
\dot m
=\frac{\dot Q}
{c_p(T_{\mathrm{return}}-T_{\mathrm{supply}})}.
\label{eq:mdot}
\end{equation}

Mass flow and volumetric flow are related by

\begin{equation}
\dot V=\frac{\dot m}{\rho}.
\end{equation}

For \(N_r\) identical racks,

\begin{equation}
\dot m_{\mathrm{total}}
=N_r\dot m_{\mathrm{rack}}.
\end{equation}

The design flow should include a margin \(f_m\):

\begin{equation}
\dot m_{\mathrm{design}}
=(1+f_m)\dot m_{\mathrm{total}}.
\end{equation}

\section{Heat-Exchanger Model}

The initial CDU and dry-cooler models use the conductance approximation

\begin{equation}
\dot Q_{\mathrm{HX}}=UA\Delta T_{\mathrm{effective}}.
\end{equation}

If inlet and outlet temperatures are available, the log-mean temperature difference
may be used:

\begin{equation}
\dot Q_{\mathrm{HX}}=UA\Delta T_{\mathrm{lm}},
\end{equation}

where

\begin{equation}
\Delta T_{\mathrm{lm}}
=\frac{\Delta T_1-\Delta T_2}
{\ln\left(\Delta T_1/\Delta T_2\right)}.
\end{equation}

For a counterflow heat exchanger,

\begin{align}
\Delta T_1&=T_{h,\mathrm{in}}-T_{c,\mathrm{out}},\\
\Delta T_2&=T_{h,\mathrm{out}}-T_{c,\mathrm{in}}.
\end{align}

When detailed manufacturer data are available, an effectiveness-NTU formulation may
replace the simple \(UA\) approximation:

\begin{equation}
\dot Q_{\mathrm{HX}}=\varepsilon C_{\min}
(T_{h,\mathrm{in}}-T_{c,\mathrm{in}}),
\end{equation}

with

\begin{equation}
C_{\min}=\min(\dot m_h c_{p,h},\dot m_c c_{p,c}),
\qquad
\mathrm{NTU}=\frac{UA}{C_{\min}}.
\end{equation}

\section{Pipe Pressure Loss}

The average coolant velocity in a circular pipe is

\begin{equation}
v=\frac{\dot V}{A_p}
=\frac{4\dot V}{\pi D^2}.
\end{equation}

The Reynolds number is

\begin{equation}
\mathrm{Re}=\frac{\rho vD}{\mu},
\end{equation}

where \(\mu\) is the dynamic viscosity. Straight-pipe pressure loss is estimated with
the Darcy--Weisbach equation:

\begin{equation}
\Delta p_{\mathrm{major}}
=f\frac{L}{D}\frac{\rho v^2}{2}.
\label{eq:darcy}
\end{equation}

Minor losses from fittings, valves, strainers, manifolds, and entrances are

\begin{equation}
\Delta p_{\mathrm{minor}}
=\sum_i K_i\frac{\rho v^2}{2}.
\end{equation}

The total loop pressure loss is

\begin{equation}
\Delta p_{\mathrm{total}}
=\Delta p_{\mathrm{major}}
+\Delta p_{\mathrm{minor}}
+\Delta p_{\mathrm{equipment}}.
\end{equation}

\section{Pump Power}

The hydraulic power transferred to the coolant is

\begin{equation}
P_{\mathrm{hydraulic}}=\Delta p_{\mathrm{total}}\dot V.
\end{equation}

Accounting for pump and motor efficiency,

\begin{equation}
P_{\mathrm{pump}}
=\frac{\Delta p_{\mathrm{total}}\dot V}
{\eta_{\mathrm{pump}}\eta_{\mathrm{motor}}}.
\label{eq:pump}
\end{equation}

For geometrically similar operation, the pump affinity laws give

\begin{align}
\frac{\dot V_2}{\dot V_1}&=\frac{N_2}{N_1},\\
\frac{H_2}{H_1}&=\left(\frac{N_2}{N_1}\right)^2,\\
\frac{P_2}{P_1}&=\left(\frac{N_2}{N_1}\right)^3,
\end{align}

where \(N\) is pump rotational speed and \(H\) is pump head.

\section{Conduction and Insulation}

One-dimensional conduction through a flat layer is modeled by Fourier's Law:

\begin{equation}
\dot Q_{\mathrm{cond}}
=\frac{kA}{L}(T_1-T_2).
\end{equation}

For radial conduction through cylindrical pipe insulation,

\begin{equation}
\dot Q_{\mathrm{pipe}}
=\frac{2\pi kL_p(T_i-T_o)}
{\ln(r_o/r_i)}.
\end{equation}

An overall thermal resistance network may be written as

\begin{equation}
\dot Q=\frac{T_{\mathrm{hot}}-T_{\mathrm{cold}}}
{R_{\mathrm{total}}},
\end{equation}

where series resistances add:

\begin{equation}
R_{\mathrm{total}}=R_{\mathrm{conv},i}
+R_{\mathrm{cond}}
+R_{\mathrm{conv},o}.
\end{equation}

\section{Residual Room Heat}

Direct-to-chip cooling does not necessarily capture all server heat. If the liquid
capture fraction is \(\phi\), then

\begin{align}
\dot Q_{\mathrm{liquid}}&=\phi P_{\mathrm{IT}},\\
\dot Q_{\mathrm{room}}&=(1-\phi)P_{\mathrm{IT}}.
\end{align}

The residual room-air node may be represented by

\begin{equation}
C_{\mathrm{air}}\frac{dT_{\mathrm{room}}}{dt}
=(1-\phi)P_{\mathrm{IT}}
-G_{\mathrm{room}}(T_{\mathrm{room}}-T_{\mathrm{supply,air}}).
\end{equation}

This term supports sizing of in-row or room-level air cooling for memory, power
supplies, networking equipment, and other non-liquid-cooled components.

\section{Facility Energy Metrics}

Total facility power is approximated as

\begin{equation}
P_{\mathrm{facility}}
=P_{\mathrm{IT}}
+P_{\mathrm{pumps}}
+P_{\mathrm{fans}}
+P_{\mathrm{chiller}}
+P_{\mathrm{UPS,loss}}
+P_{\mathrm{other}}.
\end{equation}

Power usage effectiveness is

\begin{equation}
\mathrm{PUE}
=\frac{P_{\mathrm{facility}}}{P_{\mathrm{IT}}}.
\end{equation}

Cooling-system coefficient of performance is

\begin{equation}
\mathrm{COP}
=\frac{\dot Q_{\mathrm{removed}}}{P_{\mathrm{cooling,electric}}}.
\end{equation}

A sealed dry-cooler loop is modeled with zero routine evaporative process-water use.
If an evaporative system is later considered, water use must be added as a separate
mass balance.

\section{RK4 Numerical Solver}

The coupled system is written in vector form as

\begin{equation}
\frac{d\mathbf{y}}{dt}=\mathbf{f}(t,\mathbf{y}),
\qquad
\mathbf{y}
=\begin{bmatrix}
T_g & T_s & T_p & T_{\mathrm{room}}
\end{bmatrix}^{\mathsf T}.
\end{equation}

For time step \(\Delta t\), classical fourth-order Runge--Kutta calculates

\begin{align}
\mathbf{k}_1&=\mathbf{f}(t_n,\mathbf{y}_n),\\
\mathbf{k}_2&=\mathbf{f}\left(t_n+\frac{\Delta t}{2},
\mathbf{y}_n+\frac{\Delta t}{2}\mathbf{k}_1\right),\\
\mathbf{k}_3&=\mathbf{f}\left(t_n+\frac{\Delta t}{2},
\mathbf{y}_n+\frac{\Delta t}{2}\mathbf{k}_2\right),\\
\mathbf{k}_4&=\mathbf{f}\left(t_n+\Delta t,
\mathbf{y}_n+\Delta t\mathbf{k}_3\right).
\end{align}

The next state is

\begin{equation}
\mathbf{y}_{n+1}
=\mathbf{y}_n
+\frac{\Delta t}{6}
\left(\mathbf{k}_1+2\mathbf{k}_2+2\mathbf{k}_3+\mathbf{k}_4\right).
\label{eq:rk4}
\end{equation}

RK4 has global truncation error of order \(\mathcal{O}(\Delta t^4)\) when the
solution is sufficiently smooth. Time-step convergence will be checked by repeating
the simulation with \(\Delta t\), \(\Delta t/2\), and \(\Delta t/4\).

\section{Control and Failure Functions}

GPU power may be represented using a piecewise load profile:

\begin{equation}
P_{\mathrm{IT}}(t)=
\begin{cases}
P_{\mathrm{idle}}, & 0\le t<t_1,\\
P_{\mathrm{full}}, & t_1\le t<t_2,\\
P_{\mathrm{reduced}}, & t\ge t_2.
\end{cases}
\end{equation}

A pump failure at time \(t_f\) may be modeled using

\begin{equation}
\dot m(t)=
\begin{cases}
\dot m_0, & t<t_f,\\
\alpha\dot m_0, & t\ge t_f,
\end{cases}
\end{equation}

where \(0\le\alpha<1\). For a complete flow loss, \(\alpha=0\). A standby pump
starting after delay \(t_d\) can be added as a third interval.

N+1 heat rejection can be modeled as

\begin{equation}
\dot Q_{\max,\mathrm{available}}
=N_{\mathrm{available}}\dot Q_{\max,\mathrm{unit}}.
\end{equation}

\section{Worked Flow-Rate Example}

Consider \(N_r=20\) racks, each producing \(\SI{50}{\kilo\watt}\). The total
design heat is

\begin{equation}
\dot Q_{\mathrm{total}}
=20(\SI{50}{\kilo\watt})
=\SI{1.00}{\mega\watt}.
\end{equation}

Using water with

\begin{equation}
c_p\approx\SI{4180}{\joule\per\kilogram\per\kelvin}
\end{equation}

and a design temperature rise of \(\Delta T=\SI{10}{\kelvin}\),

\begin{align}
\dot m
&=\frac{\SI{1.00e6}{\watt}}
{(\SI{4180}{\joule\per\kilogram\per\kelvin})(\SI{10}{\kelvin})}\\
&\approx\SI{23.9}{\kilogram\per\second}.
\end{align}

Taking \(\rho\approx\SI{1000}{\kilogram\per\cubic\metre}\),

\begin{equation}
\dot V
=\frac{23.9}{1000}
=\SI{0.0239}{\cubic\metre\per\second}
\approx\SI{379}{\gpm}.
\end{equation}

With a \(15\%\) design margin,

\begin{equation}
\dot V_{\mathrm{design}}
=1.15(\SI{379}{\gpm})
\approx\SI{436}{\gpm}.
\end{equation}

This is a conceptual sizing calculation. Final design flow depends on the selected
coolant, allowable temperature rise, equipment curves, redundancy, and pipe losses.

\section{Inputs and Outputs}

\subsection{Principal Inputs}

\begin{itemize}[leftmargin=2em]
    \item Rack count, GPU count, and load profile
    \item GPU and coolant initial temperatures
    \item Thermal capacitances \(C_g\), \(C_s\), and \(C_p\)
    \item Conductances \(G_{gs}\), \(G_{sp}\), and \(G_{pa}\)
    \item Coolant density, specific heat, and viscosity
    \item Secondary and primary flow rates
    \item Pipe lengths, diameters, roughness, and fitting loss coefficients
    \item Pump efficiency and equipment pressure losses
    \item Chiller or dry-cooler maximum heat-rejection capacity
    \item Outdoor temperature profile
    \item Maximum permitted GPU and coolant temperatures
\end{itemize}

\subsection{Principal Outputs}

\begin{itemize}[leftmargin=2em]
    \item GPU, secondary-coolant, primary-loop, and room temperatures versus time
    \item Peak temperature and time of occurrence
    \item Time before a thermal safety limit is exceeded
    \item Heat generated, transported, stored, and rejected
    \item Required rack and facility coolant flow
    \item Loop pressure loss and pump power
    \item Cooling electrical power, COP, and estimated PUE
    \item Pass or fail result for normal and failure scenarios
\end{itemize}

\section{Modeling Assumptions}

The first version of the model will use the following assumptions:

\begin{enumerate}[leftmargin=2em]
    \item Each thermal node is spatially uniform.
    \item Coolant properties are constant over the operating temperature range.
    \item There is no phase change in the secondary or primary coolant loops.
    \item Nearly all IT electrical power becomes heat.
    \item Effective conductances represent combined cold-plate, convection, contact,
          and heat-exchanger resistances.
    \item Heat-rejection equipment has a finite rated capacity.
    \item The closed loop has no routine evaporative water loss unless an evaporative
          heat-rejection system is explicitly selected.
    \item Radiation is small relative to forced convection and liquid heat transport
          and is omitted from the initial model.
\end{enumerate}

\section{Validation Plan}

\begin{enumerate}[leftmargin=2em]
    \item Compare the RK4 result for Equation~\eqref{eq:one-node} with the analytical
          solution in Equation~\eqref{eq:analytical}.
    \item Perform time-step convergence testing.
    \item Confirm energy conservation by integrating generated and rejected heat.
    \item Compare steady-state flow and heat rejection with
          Equation~\eqref{eq:flowheat}.
    \item Compare pump calculations with manufacturer performance curves.
    \item Replace placeholder conductances and capacities with manufacturer data.
    \item If possible, use a small-scale step-load experiment to estimate thermal
          capacitance and conductance.
\end{enumerate}

\section{Revit Integration}

Each rack family should store the parameters

\begin{center}
Rack ID, GPU count, electrical power, design heat load, coolant flow, supply
temperature, return temperature, and assigned CDU.
\end{center}

Revit schedules will be exported to CSV and used as model inputs. The computed flow,
capacity, and redundancy results will be used to revise pipe sizes, pump selection,
CDU quantities, equipment placement, and clearances in the Revit model.

\section{Proposed Simulation Cases}

\begin{enumerate}[leftmargin=2em]
    \item Normal steady operation
    \item Idle-to-full-load GPU step
    \item Phoenix summer outdoor-temperature condition
    \item Reduced coolant flow
    \item Failure of one pump
    \item Failure of one CDU
    \item Loss of outdoor heat rejection
    \item Emergency IT shutdown
    \item N+1 redundant cooling operation
    \item Restart after a cooling interruption
\end{enumerate}

\section{Conclusion}

The OASIS mathematical model connects the Revit design to quantitative thermal and
fluid performance. Conservation of energy, Newton's Law of Cooling, sensible heat
transport, heat-exchanger relations, Darcy--Weisbach pressure loss, pump power, and
RK4 integration provide a defensible framework for evaluating a scaled GPU cooling
system. The model will show not only the expected operating temperature, but also
how rapidly the system responds to computational load changes and how long it remains
safe after a cooling-system failure.

\end{document}
